\documentclass[11pt]{amsart}

\usepackage[T1]{fontenc}
\usepackage{lmodern}
\usepackage{microtype}
\usepackage{amsmath,amssymb,amsthm}
\usepackage[hidelinks]{hyperref}

\hypersetup{
  pdftitle={A Counterexample for [65] to an Exhaustiveness Conjecture
    of Espinosa-Garcia and Pellicer},
  pdfkeywords={Sidon set, Sidon--Ramsey partition, balanced partition,
    counterexample}
}

\newtheorem{theorem}{Theorem}

\title[A counterexample to an exhaustiveness conjecture]
{A Counterexample for $[65]$ to an Exhaustiveness Conjecture
of Espinosa-Garc\'ia and Pellicer}

\subjclass[2020]{Primary 11B75; Secondary 05D10}
\keywords{Sidon set, Sidon--Ramsey partition, balanced partition, counterexample}
\date{}

\begin{document}

\begin{abstract}
We exhibit a balanced partition of $[65]$ into seven Sidon sets that is
not among the five partitions displayed by Espinosa-Garc\'ia and Pellicer
or their reflections. This refutes their exhaustiveness conjecture in
arXiv:2309.08553v1; reflection gives a second distinct omitted partition.
\end{abstract}

\maketitle

Write $[n]=\{1,\ldots,n\}$. A finite set $A\subset\mathbb Z$ is
\emph{Sidon} if every integer has at most one representation $a+b$ with
$a,b\in A$ and $a\le b$. A partition of $[n]$ into Sidon sets is a
\emph{Sidon--Ramsey partition}; it is \emph{balanced} if any two part
sizes differ by at most one.

Espinosa-Garc\'ia and Pellicer display five balanced partitions of
$[65]$ into seven Sidon sets and then write, in unnumbered prose
immediately after the displays \cite[arXiv v1, p.~6]{EGP}:
\begin{quote}
These partitions have two parts of $10$ elements each, and five parts
with $9$ elements each. We conjecture that these and their reflected
partitions (constructed by including the numbers $66-x$ instead of $x$
in each part) are all balanced Sidon-Ramsey partitions with those
parameters.
\end{quote}
We call this their \emph{exhaustiveness conjecture}; it carries no
number in \cite{EGP}. Since $65=7\cdot9+2$, every such partition has two
parts of size $10$ and five of size $9$, in agreement with the quoted
sentence.

\begin{theorem}
The exhaustiveness conjecture quoted above is false.
\end{theorem}

\begin{proof}
Let $\mathcal P=\{A_1,\ldots,A_7\}$, where
\begin{align*}
A_1&=\{1,5,13,23,26,28,37,56,57,63\},&
A_2&=\{2,20,21,24,33,44,50,58,60,65\},\\
A_3&=\{3,4,27,30,32,39,43,49,64\},&
A_4&=\{6,15,16,22,36,40,48,51,53\},\\
A_5&=\{7,14,19,25,29,38,52,54,55\},&
A_6&=\{8,11,17,34,41,45,46,59,61\},\\
A_7&=\{9,10,12,18,31,35,42,47,62\}.&&
\end{align*}
The listed integers are pairwise distinct and have union $[65]$; the
part sizes are $10,10,9,9,9,9,9$.

For finite $A\subset\mathbb Z$, put
\[
  \Sigma(A)=\{a+b:a,b\in A,\ a\le b\}.
\]
Then $A$ is Sidon exactly when
$|\Sigma(A)|=\binom{|A|+1}{2}$. Exact addition gives
\[
\begin{array}{c|ccccccc}
 i&1&2&3&4&5&6&7\\ \hline
 |A_i|&10&10&9&9&9&9&9\\
 |\Sigma(A_i)|&55&55&45&45&45&45&45
\end{array}
\]
so every $A_i$ is Sidon and $\mathcal P$ is a balanced Sidon--Ramsey
partition of $[65]$. Equivalently, $A$ is Sidon precisely when the
$\binom{|A|}{2}$ differences $b-a$ with $a<b$ in $A$ are pairwise
distinct. That is a check on $45$ differences for each of $A_1$ and
$A_2$ and on $36$ for each of the remaining five parts, $270$ in all,
and in each part these differences are pairwise distinct.

We regard partitions as unordered families of parts. The ten
$10$-element parts occurring in the five partitions displayed in
\cite[arXiv v1, p.~6]{EGP} are the eight distinct sets
\begin{align*}
 &\{1,3,15,22,30,33,46,50,55,56\},&
&\{2,4,16,23,31,34,47,51,56,57\},\\
 &\{2,9,17,21,26,27,47,49,60,63\},&
&\{2,6,14,24,27,29,38,57,58,64\},\\
 &\{7,10,15,19,33,43,44,50,63,65\},&
&\{3,9,10,29,38,40,43,53,61,65\},\\
 &\{3,5,17,24,32,35,48,52,57,58\},&
&\{2,6,14,16,19,37,38,44,53,64\},
\end{align*}
of which the first two occur twice: the first is a part of the first and
of the second displayed partition, the second a part of the third and of
the fourth. None of the eight contains $\{2,65\}$ and none contains
$\{1,64\}$. Therefore no $10$-element part in those partitions or their
reflections contains $\{2,65\}$. Since $\{2,65\}\subset A_2$, the
partition $\mathcal P$ is not in the conjectured family.

Let $\rho(x)=66-x$ and apply $\rho$ elementwise to sets and partitions.
Reflection preserves sizes and the Sidon property, and the conjectured
family is $\rho$-invariant. Hence $\rho(\mathcal P)$ is another omitted
partition. It is distinct from $\mathcal P$: the $10$-element part
$\rho(A_2)$ contains $\{1,64\}$, whereas neither $10$-element part of
$\mathcal P$ does.
\end{proof}

The result concerns precisely the statement quoted above, as it appears
on p.~6 of arXiv:2309.08553v1; no complete census of the balanced
Sidon--Ramsey partitions of $[65]$ with seven parts is asserted.

We have not been able to consult the published version listed in
\cite{EGP}, and the e-print stands unrevised at v1, so whether the
quoted sentence appears unchanged there is left open. The mathematical
content does not depend on this: $\mathcal P$ and $\rho(\mathcal P)$ are
balanced Sidon--Ramsey partitions of $[65]$ into seven Sidon sets, and
by the argument above neither is one of the five partitions displayed in
\cite{EGP} nor a reflection of one.

\begin{thebibliography}{1}

\bibitem{EGP}
M.~A. Espinosa-Garc\'ia and D.~Pellicer,
\emph{Update on Sidon--Ramsey numbers},
\href{https://arxiv.org/abs/2309.08553v1}{arXiv:2309.08553v1}
[math.CO] (2023); published in
\emph{Discrete Appl. Math.} \textbf{378} (2026), 120--124,
\href{https://doi.org/10.1016/j.dam.2025.07.002}
{doi:10.1016/j.dam.2025.07.002}.

\end{thebibliography}

\end{document}